\documentclass{article}
\usepackage[utf8]{inputenc}
\usepackage[T1]{fontenc}
\usepackage{amsmath, amssymb, amsthm}
\usepackage{geometry}
\usepackage{graphicx}
\usepackage{fancyhdr}
\usepackage{xcolor}
\usepackage{enumitem}
\usepackage{cancel}
\usepackage{hyperref}

% Geometry setup
\geometry{a4paper, margin=1in}

% Header and Footer
\pagestyle{fancy}
\fancyhf{}
\lhead{Homework 14.6: Directional Derivatives and Gradient}
\rhead{Tashfeen Omran}
\rfoot{Page \thepage}

% Title Information
\title{\textbf{Homework 14.6: Directional Derivatives and Gradient}\\ \large A Comprehensive Study Guide and Solution Manual}
\author{Tashfeen Omran}
\date{February 2026}

\begin{document}

\maketitle
\tableofcontents
\newpage

\section{Comprehensive Introduction: Context and Prerequisites}

\subsection{Core Concepts: Beyond Partial Derivatives}
In previous chapters, we learned how to calculate partial derivatives, such as $f_x$ and $f_y$. These measure the rate of change of a function $f(x, y)$ if we move purely in the $x$-direction (East) or purely in the $y$-direction (North).

But what if we want to move Northeast? Or at an angle of $30^\circ$?
The \textbf{Directional Derivative}, denoted $D_{\mathbf{u}}f$, generalizes the concept of a derivative to \textit{any} arbitrary direction defined by a unit vector $\mathbf{u}$.

To calculate this efficiently, we introduce the \textbf{Gradient Vector}, denoted $\nabla f$ (read "del f"). The gradient bundles all the partial derivatives into a single vector. It is the fundamental tool of multivariable calculus, acting as a "compass" that always points in the direction of steepest ascent.

\subsection{Intuition and Derivation}
Imagine you are standing on a hill defined by $z = f(x, y)$.
\begin{itemize}
    \item $f_x$ is the slope if you look East.
    \item $f_y$ is the slope if you look North.
    \item If you face a direction given by a unit vector $\mathbf{u} = \langle a, b \rangle$, the slope you see is a "weighted average" of the slopes in the $x$ and $y$ directions, weighted by how much your direction aligns with them.
\end{itemize}

Mathematically, if $\mathbf{u} = \langle a, b \rangle$ is a unit vector ($a^2 + b^2 = 1$), the directional derivative is:
\[ D_{\mathbf{u}}f(x, y) = f_x(x, y)a + f_y(x, y)b \]
This is exactly the dot product of the gradient vector $\nabla f = \langle f_x, f_y \rangle$ and the direction vector $\mathbf{u}$.

\subsection{Historical Context and Motivation}
In the 19th century, mathematicians like Hamilton and Gibbs needed tools to describe physical fields—like gravity, fluid flow, and electromagnetism—where forces act in 3D space, not just along axes. They needed a way to describe "change" independent of the coordinate system. The gradient vector was developed to capture the intrinsic geometry of a scalar field, allowing physicists to calculate potential energy changes or heat flow in any direction without being tied to a specific grid.

\subsection{Key Formulas}
\begin{enumerate}
    \item \textbf{The Gradient Vector (2D and 3D):}
    \[ \nabla f(x, y) = \langle f_x, f_y \rangle = \frac{\partial f}{\partial x}\mathbf{i} + \frac{\partial f}{\partial y}\mathbf{j} \]
    \[ \nabla f(x, y, z) = \langle f_x, f_y, f_z \rangle \]

    \item \textbf{Directional Derivative (Dot Product Form):}
    \[ D_{\mathbf{u}}f = \nabla f \cdot \mathbf{u} \]
    \textbf{Crucial:} The vector $\mathbf{u}$ \textbf{MUST} be a unit vector ($|\mathbf{u}|=1$). If you are given a general vector $\mathbf{v}$, you must normalize it: $\mathbf{u} = \frac{\mathbf{v}}{|\mathbf{v}|}$.

    \item \textbf{Directional Derivative (Angle Form):}
    If the direction is given by an angle $\theta$ (counterclockwise from positive $x$-axis):
    \[ \mathbf{u} = \langle \cos\theta, \sin\theta \rangle \implies D_{\mathbf{u}}f = f_x \cos\theta + f_y \sin\theta \]

    \item \textbf{Maximizing Rate of Change:}
    \begin{itemize}
        \item \textbf{Max Rate:} $|\nabla f|$ (The magnitude of the gradient).
        \item \textbf{Direction of Max Rate:} The direction of $\nabla f$.
    \end{itemize}

    \item \textbf{Tangent Planes to Level Surfaces:}
    For a surface defined implicitly by $F(x, y, z) = k$, the normal vector is $\mathbf{n} = \nabla F$.
    The tangent plane at $(x_0, y_0, z_0)$ is:
    \[ F_x(x_0, y_0, z_0)(x - x_0) + F_y(x_0, y_0, z_0)(y - y_0) + F_z(x_0, y_0, z_0)(z - z_0) = 0 \]
\end{enumerate}

\subsection{Prerequisites}
\begin{itemize}
    \item \textbf{Partial Differentiation:} Ability to compute $f_x, f_y, f_z$ quickly (Chain Rule, Product Rule).
    \item \textbf{Vector Operations:} Dot product, magnitude, and unit vector normalization.
    \item \textbf{Trigonometry:} Values of $\sin$ and $\cos$ for standard unit circle angles.
\end{itemize}

\newpage

\section{Detailed Homework Solutions}

\subsection*{Problem 1}
\textbf{Find the directional derivative of $f(x, y) = xy^3 - x^2$ at $(1, 4)$ in the direction $\theta = \frac{\pi}{3}$.}

\textbf{Step 1: Find the Gradient $\nabla f$.}
\[ f_x = \frac{\partial}{\partial x}(xy^3 - x^2) = y^3 - 2x \]
\[ f_y = \frac{\partial}{\partial y}(xy^3 - x^2) = 3xy^2 \]
\[ \nabla f(x, y) = \langle y^3 - 2x, 3xy^2 \rangle \]

\textbf{Step 2: Evaluate Gradient at $(1, 4)$.}
\[ \nabla f(1, 4) = \langle 4^3 - 2(1), 3(1)(4^2) \rangle = \langle 64 - 2, 3(16) \rangle = \langle 62, 48 \rangle \]

\textbf{Step 3: Determine Unit Vector $\mathbf{u}$.}
Given $\theta = \frac{\pi}{3}$:
\[ \mathbf{u} = \langle \cos(\pi/3), \sin(\pi/3) \rangle = \left\langle \frac{1}{2}, \frac{\sqrt{3}}{2} \right\rangle \]

\textbf{Step 4: Calculate $D_{\mathbf{u}}f = \nabla f \cdot \mathbf{u}$.}
\[ D_{\mathbf{u}}f(1, 4) = 62\left(\frac{1}{2}\right) + 48\left(\frac{\sqrt{3}}{2}\right) = 31 + 24\sqrt{3} \]

\textbf{Final Answer:}
\[ \boxed{31 + 24\sqrt{3}} \]

\hrulefill

\subsection*{Problem 2}
\textbf{Consider $f(x, y) = xy^3 - x^2$. Find partial derivatives, and $D_{\mathbf{u}}f$ at $(1, 2)$ for $\theta = \frac{\pi}{3}$.}

\textbf{Part (a): Find partial derivatives.}
\[ f_x(x, y) = y^3 - 2x \]
\[ f_y(x, y) = 3xy^2 \]

\textbf{Part (b): Find $D_{\mathbf{u}}f(1, 2)$ with $\theta = \frac{\pi}{3}$.}
Evaluate gradient at $(1, 2)$:
\[ \nabla f(1, 2) = \langle 2^3 - 2(1), 3(1)(2^2) \rangle = \langle 8 - 2, 12 \rangle = \langle 6, 12 \rangle \]
Unit vector $\mathbf{u} = \langle \frac{1}{2}, \frac{\sqrt{3}}{2} \rangle$.
\[ D_{\mathbf{u}}f = \langle 6, 12 \rangle \cdot \left\langle \frac{1}{2}, \frac{\sqrt{3}}{2} \right\rangle = 3 + 6\sqrt{3} \]

\textbf{Final Answer:}
\[ f_x = \boxed{y^3 - 2x}, \quad f_y = \boxed{3xy^2}, \quad D_{\mathbf{u}}f(1, 2) = \boxed{3 + 6\sqrt{3}} \]

\hrulefill

\subsection*{Problem 3}
\textbf{Find $D_{\mathbf{u}}f$ for $f(x, y) = y \cos(xy)$ at $(0, 1)$ with $\theta = \frac{\pi}{4}$.}

\textbf{Step 1: Gradient.}
\[ f_x = y(-\sin(xy) \cdot y) = -y^2 \sin(xy) \]
\[ f_y = 1 \cdot \cos(xy) + y(-\sin(xy) \cdot x) = \cos(xy) - xy \sin(xy) \]

\textbf{Step 2: Evaluate at $(0, 1)$.}
\[ f_x(0, 1) = -(1)^2 \sin(0) = 0 \]
\[ f_y(0, 1) = \cos(0) - 0 = 1 \]
\[ \nabla f(0, 1) = \langle 0, 1 \rangle \]

\textbf{Step 3: Direction.}
\[ \mathbf{u} = \langle \cos(\pi/4), \sin(\pi/4) \rangle = \left\langle \frac{\sqrt{2}}{2}, \frac{\sqrt{2}}{2} \right\rangle \]

\textbf{Step 4: Dot Product.}
\[ D_{\mathbf{u}}f = 0\left(\frac{\sqrt{2}}{2}\right) + 1\left(\frac{\sqrt{2}}{2}\right) = \frac{\sqrt{2}}{2} \]

\textbf{Final Answer:}
\[ \boxed{\frac{\sqrt{2}}{2}} \]

\hrulefill

\subsection*{Problem 4}
\textbf{Find $D_{\mathbf{u}}f$ for $f(x, y) = \sqrt{8x + 9y}$ at $(9, 1)$ with $\theta = -\frac{\pi}{6}$.}

\textbf{Step 1: Gradient.}
\[ f_x = \frac{1}{2\sqrt{8x+9y}} \cdot 8 = \frac{4}{\sqrt{8x+9y}} \]
\[ f_y = \frac{1}{2\sqrt{8x+9y}} \cdot 9 = \frac{4.5}{\sqrt{8x+9y}} \]

\textbf{Step 2: Evaluate at $(9, 1)$.}
Term inside root: $8(9) + 9(1) = 72 + 9 = 81$. $\sqrt{81} = 9$.
\[ f_x(9, 1) = \frac{4}{9}, \quad f_y(9, 1) = \frac{4.5}{9} = \frac{1}{2} \]
\[ \nabla f = \langle \frac{4}{9}, \frac{1}{2} \rangle \]

\textbf{Step 3: Direction.}
\[ \mathbf{u} = \langle \cos(-\pi/6), \sin(-\pi/6) \rangle = \left\langle \frac{\sqrt{3}}{2}, -\frac{1}{2} \right\rangle \]

\textbf{Step 4: Dot Product.}
\[ D_{\mathbf{u}}f = \left(\frac{4}{9}\right)\left(\frac{\sqrt{3}}{2}\right) + \left(\frac{1}{2}\right)\left(-\frac{1}{2}\right) = \frac{2\sqrt{3}}{9} - \frac{1}{4} \]

\textbf{Final Answer:}
\[ \boxed{\frac{2\sqrt{3}}{9} - \frac{1}{4}} \]

\hrulefill

\subsection*{Problem 5}
\textbf{Find $D_{\mathbf{u}}f$ for $f(x, y) = \arctan(xy)$ at $(3, -5)$ with $\theta = \frac{3\pi}{4}$.}

\textbf{Step 1: Gradient.}
Recall $\frac{d}{du}\arctan(u) = \frac{1}{1+u^2}$.
\[ f_x = \frac{1}{1+(xy)^2}(y), \quad f_y = \frac{1}{1+(xy)^2}(x) \]

\textbf{Step 2: Evaluate at $(3, -5)$.}
$xy = -15$, $(xy)^2 = 225$. Denom = $226$.
\[ f_x = \frac{-5}{226}, \quad f_y = \frac{3}{226} \]

\textbf{Step 3: Direction.}
\[ \mathbf{u} = \langle \cos(3\pi/4), \sin(3\pi/4) \rangle = \left\langle -\frac{\sqrt{2}}{2}, \frac{\sqrt{2}}{2} \right\rangle \]

\textbf{Step 4: Dot Product.}
\[ D_{\mathbf{u}}f = \left(\frac{-5}{226}\right)\left(\frac{-\sqrt{2}}{2}\right) + \left(\frac{3}{226}\right)\left(\frac{\sqrt{2}}{2}\right) = \frac{5\sqrt{2}}{452} + \frac{3\sqrt{2}}{452} = \frac{8\sqrt{2}}{452} \]
Simplify fraction: $8/452 = 2/113$.

\textbf{Final Answer:}
\[ \boxed{\frac{2\sqrt{2}}{113}} \]

\hrulefill

\subsection*{Problem 6}
\textbf{$f(x, y) = x^2 e^y$, $P(7, 0)$, $\mathbf{u} = \frac{1}{5}(3\mathbf{i} - 4\mathbf{j})$.}

\textbf{(a) Find $\nabla f(x, y)$.}
\[ f_x = 2x e^y, \quad f_y = x^2 e^y \]
\[ \nabla f = \langle 2x e^y, x^2 e^y \rangle \]

\textbf{(b) Evaluate at $(7, 0)$.}
$e^0 = 1$.
\[ \nabla f(7, 0) = \langle 2(7)(1), 7^2(1) \rangle = \langle 14, 49 \rangle \]

\textbf{(c) Rate of change ($D_{\mathbf{u}}f$).}
The given vector $\mathbf{u} = \langle \frac{3}{5}, -\frac{4}{5} \rangle$ is already a unit vector ($|\mathbf{u}| = \sqrt{(3/5)^2 + (-4/5)^2} = 1$).
\[ D_{\mathbf{u}}f = \langle 14, 49 \rangle \cdot \langle \frac{3}{5}, -\frac{4}{5} \rangle = \frac{42}{5} - \frac{196}{5} = \frac{-154}{5} \]

\textbf{Final Answer:}
(a) $\langle 2xe^y, x^2e^y \rangle$ \quad (b) $\langle 14, 49 \rangle$ \quad (c) $\boxed{-30.8}$ or $\boxed{-\frac{154}{5}}$

\hrulefill

\subsection*{Problem 7}
\textbf{$f(x, y) = x/y$, $P(3, 1)$, $\mathbf{u} = \frac{3}{5}\mathbf{i} + \frac{4}{5}\mathbf{j}$.}

\textbf{(a) $\nabla f(x, y)$.}
\[ f_x = \frac{1}{y}, \quad f_y = x(-y^{-2}) = -\frac{x}{y^2} \]
\[ \nabla f = \langle \frac{1}{y}, -\frac{x}{y^2} \rangle \]

\textbf{(b) $\nabla f(3, 1)$.}
\[ \nabla f(3, 1) = \langle \frac{1}{1}, -\frac{3}{1^2} \rangle = \langle 1, -3 \rangle \]

\textbf{(c) $D_{\mathbf{u}}f$.}
$\mathbf{u} = \langle \frac{3}{5}, \frac{4}{5} \rangle$ is a unit vector.
\[ D_{\mathbf{u}}f = 1(\frac{3}{5}) + (-3)(\frac{4}{5}) = \frac{3}{5} - \frac{12}{5} = -\frac{9}{5} \]

\textbf{Final Answer:}
(a) $\langle \frac{1}{y}, -\frac{x}{y^2} \rangle$ \quad (b) $\langle 1, -3 \rangle$ \quad (c) $\boxed{-\frac{9}{5}}$

\hrulefill

\subsection*{Problem 8}
\textbf{$f(x, y, z) = x^2 yz - xyz^8$, $P(6, -1, 1)$, $\mathbf{u} = \langle 0, \frac{4}{5}, -\frac{3}{5} \rangle$.}

\textbf{(a) $\nabla f$.}
\[ f_x = 2xyz - yz^8 \]
\[ f_y = x^2z - xz^8 \]
\[ f_z = x^2y - 8xyz^7 \]

\textbf{(b) $\nabla f(6, -1, 1)$.}
\[ f_x = 2(6)(-1)(1) - (-1)(1) = -12 + 1 = -11 \]
\[ f_y = 36(1) - 6(1) = 30 \]
\[ f_z = 36(-1) - 8(6)(-1)(1) = -36 + 48 = 12 \]
\[ \nabla f = \langle -11, 30, 12 \rangle \]

\textbf{(c) $D_{\mathbf{u}}f$.}
$\mathbf{u}$ is a unit vector.
\[ D_{\mathbf{u}}f = -11(0) + 30\left(\frac{4}{5}\right) + 12\left(-\frac{3}{5}\right) = 0 + 24 - \frac{36}{5} = \frac{120 - 36}{5} = \frac{84}{5} \]

\textbf{Final Answer:}
(b) $\langle -11, 30, 12 \rangle$ \quad (c) $\boxed{\frac{84}{5}}$

\hrulefill

\subsection*{Problem 9}
\textbf{$f(x, y, z) = y^2 e^{xyz}$, $P(0, 3, -3)$, $\mathbf{u} = \langle \frac{4}{13}, \frac{12}{13}, \frac{3}{13} \rangle$.}

\textbf{(a) $\nabla f$.}
\[ f_x = y^2 e^{xyz}(yz) = y^3 z e^{xyz} \]
\[ f_y = 2y e^{xyz} + y^2 e^{xyz}(xz) = (2y + xy^2z)e^{xyz} \]
\[ f_z = y^2 e^{xyz}(xy) = xy^3 e^{xyz} \]

\textbf{(b) $\nabla f(0, 3, -3)$.}
Exponent $xyz = 0(3)(-3) = 0$. $e^0 = 1$.
\[ f_x = 3^3(-3)(1) = 27(-3) = -81 \]
\[ f_y = (2(3) + 0) = 6 \]
\[ f_z = 0(27)(1) = 0 \]
\[ \nabla f = \langle -81, 6, 0 \rangle \]

\textbf{(c) $D_{\mathbf{u}}f$.}
\[ D_{\mathbf{u}}f = -81\left(\frac{4}{13}\right) + 6\left(\frac{12}{13}\right) + 0 = \frac{-324 + 72}{13} = \frac{-252}{13} \]

\textbf{Final Answer:}
(b) $\langle -81, 6, 0 \rangle$ \quad (c) $\boxed{-\frac{252}{13}}$

\hrulefill

\subsection*{Problem 10}
\textbf{$f(x, y) = 3e^x \sin(y)$, $(0, \pi/3)$, $\mathbf{v} = \langle -3, 4 \rangle$.}

\textbf{Caution:} $\mathbf{v}$ is NOT a unit vector.
$|\mathbf{v}| = \sqrt{(-3)^2 + 4^2} = \sqrt{9+16} = 5$.
$\mathbf{u} = \langle -3/5, 4/5 \rangle$.

\textbf{Step 1: Gradient.}
\[ \nabla f = \langle 3e^x \sin y, 3e^x \cos y \rangle \]
At $(0, \pi/3)$: $e^0 = 1$, $\sin(\pi/3)=\frac{\sqrt{3}}{2}$, $\cos(\pi/3)=\frac{1}{2}$.
\[ \nabla f = \langle \frac{3\sqrt{3}}{2}, \frac{3}{2} \rangle \]

\textbf{Step 2: $D_{\mathbf{u}}f$.}
\[ D_{\mathbf{u}}f = \left(\frac{3\sqrt{3}}{2}\right)\left(-\frac{3}{5}\right) + \left(\frac{3}{2}\right)\left(\frac{4}{5}\right) = \frac{-9\sqrt{3}}{10} + \frac{12}{10} = \frac{12 - 9\sqrt{3}}{10} \]

\textbf{Final Answer:}
\[ \boxed{\frac{12 - 9\sqrt{3}}{10}} \]

\hrulefill

\subsection*{Problem 11}
\textbf{$f(x, y) = 7e^x \sin(y)$, $(0, \pi/3)$, $\mathbf{v} = \langle -10, 24 \rangle$.}

1. \textbf{Gradient $\nabla f(x,y)$:} $\langle 7e^x \sin y, 7e^x \cos y \rangle$.
2. \textbf{Eval at $(0, \pi/3)$:} $\langle 7\frac{\sqrt{3}}{2}, 7\frac{1}{2} \rangle = \langle \frac{7\sqrt{3}}{2}, \frac{7}{2} \rangle$.
3. \textbf{Unit Vector $\mathbf{u}$:} $|\mathbf{v}| = \sqrt{100 + 576} = \sqrt{676} = 26$.
   $\mathbf{u} = \frac{1}{26}\langle -10, 24 \rangle = \langle -\frac{5}{13}, \frac{12}{13} \rangle$.
4. \textbf{$D_{\mathbf{u}}f$:}
   \[ \left(\frac{7\sqrt{3}}{2}\right)\left(-\frac{5}{13}\right) + \left(\frac{7}{2}\right)\left(\frac{12}{13}\right) = \frac{-35\sqrt{3}}{26} + \frac{84}{26} = \frac{84 - 35\sqrt{3}}{26} \]

\textbf{Final Answer:}
$\nabla f = \langle 7e^x \sin y, 7e^x \cos y \rangle$, $\nabla f(P) = \langle 3.5\sqrt{3}, 3.5 \rangle$, $\mathbf{u} = \langle -5/13, 12/13 \rangle$, $D_{\mathbf{u}}f = \boxed{\frac{84 - 35\sqrt{3}}{26}}$

\hrulefill

\subsection*{Problem 12}
\textbf{$f(x, y) = \frac{x}{x^2+y^2}$, $(2, 1)$, $\mathbf{v} = \langle 5, 4 \rangle$.}

\textbf{Step 1: Gradient.} Quotient Rule.
\[ f_x = \frac{1(x^2+y^2) - x(2x)}{(x^2+y^2)^2} = \frac{y^2 - x^2}{(x^2+y^2)^2} \]
\[ f_y = \frac{0 - x(2y)}{(x^2+y^2)^2} = \frac{-2xy}{(x^2+y^2)^2} \]

\textbf{Step 2: Eval at $(2, 1)$.} Denom $(2^2+1^2)^2 = 25$.
\[ f_x = \frac{1 - 4}{25} = \frac{-3}{25} \]
\[ f_y = \frac{-4}{25} \]

\textbf{Step 3: Unit Vector.}
$|\mathbf{v}| = \sqrt{25+16} = \sqrt{41}$.
$\mathbf{u} = \langle \frac{5}{\sqrt{41}}, \frac{4}{\sqrt{41}} \rangle$.

\textbf{Step 4: Dot Product.}
\[ D_{\mathbf{u}}f = \left(\frac{-3}{25}\right)\left(\frac{5}{\sqrt{41}}\right) + \left(\frac{-4}{25}\right)\left(\frac{4}{\sqrt{41}}\right) = \frac{-15 - 16}{25\sqrt{41}} = \frac{-31}{25\sqrt{41}} \]

\textbf{Final Answer:}
\[ \boxed{-\frac{31}{25\sqrt{41}}} \]

\hrulefill

\subsection*{Problem 13}
\textbf{$g(s, t) = s\sqrt{t}$, $(3, 9)$, $\mathbf{v} = 3\mathbf{i} - \mathbf{j}$.}

\textbf{Step 1: Gradient.}
\[ g_s = \sqrt{t}, \quad g_t = s \cdot \frac{1}{2\sqrt{t}} \]
At $(3, 9)$: $\sqrt{9}=3$.
\[ g_s = 3, \quad g_t = \frac{3}{6} = \frac{1}{2} \]
$\nabla g = \langle 3, 0.5 \rangle$.

\textbf{Step 2: Unit Vector.}
$|\mathbf{v}| = \sqrt{9+1} = \sqrt{10}$. $\mathbf{u} = \langle \frac{3}{\sqrt{10}}, \frac{-1}{\sqrt{10}} \rangle$.

\textbf{Step 3: Result.}
\[ D_{\mathbf{u}}g = 3\left(\frac{3}{\sqrt{10}}\right) + 0.5\left(\frac{-1}{\sqrt{10}}\right) = \frac{9 - 0.5}{\sqrt{10}} = \frac{8.5}{\sqrt{10}} = \frac{17}{2\sqrt{10}} \]

\textbf{Final Answer:}
\[ \boxed{\frac{17}{2\sqrt{10}}} \]

\hrulefill

\subsection*{Problem 14}
\textbf{$h(r, s, t) = \ln(3r + 6s + 9t)$, $(2, 2, 2)$, $\mathbf{v} = 16\mathbf{i} + 48\mathbf{j} + 24\mathbf{k}$.}

\textbf{Step 1: Gradient.}
Inside: $3(2)+6(2)+9(2) = 6+12+18 = 36$.
\[ h_r = \frac{3}{36} = \frac{1}{12}, \quad h_s = \frac{6}{36} = \frac{1}{6}, \quad h_t = \frac{9}{36} = \frac{1}{4} \]
$\nabla h = \langle \frac{1}{12}, \frac{1}{6}, \frac{1}{4} \rangle$.

\textbf{Step 2: Unit Vector.}
$\mathbf{v} = 8\langle 2, 6, 3 \rangle$. Magnitude is $8\sqrt{4+36+9} = 8\sqrt{49} = 8(7) = 56$.
$\mathbf{u} = \frac{1}{56}\langle 16, 48, 24 \rangle = \langle \frac{2}{7}, \frac{6}{7}, \frac{3}{7} \rangle$.

\textbf{Step 3: Dot Product.}
\[ \frac{1}{12}(\frac{2}{7}) + \frac{1}{6}(\frac{6}{7}) + \frac{1}{4}(\frac{3}{7}) = \frac{1}{42} + \frac{1}{7} + \frac{3}{28} \]
Common denom 84:
\[ \frac{2}{84} + \frac{12}{84} + \frac{9}{84} = \frac{23}{84} \]

\textbf{Final Answer:}
\[ \boxed{\frac{23}{84}} \]

\hrulefill

\subsection*{Problem 15 \& 16}
\textbf{$f(x, y) = x/y^2$, $P(3, -1)$, toward $Q(-6, 11)$.}

1. \textbf{Gradient:} $f_x = y^{-2}$, $f_y = -2xy^{-3}$.
   At $(3, -1)$: $f_x = (-1)^{-2} = 1$, $f_y = -2(3)(-1)^{-3} = -6(-1) = 6$.
   $\nabla f = \langle 1, 6 \rangle$.

2. \textbf{Vector $\vec{PQ}$:} $\langle -6-3, 11-(-1) \rangle = \langle -9, 12 \rangle$.
   Magnitude: $\sqrt{81 + 144} = \sqrt{225} = 15$.
   $\mathbf{u} = \langle \frac{-9}{15}, \frac{12}{15} \rangle = \langle -\frac{3}{5}, \frac{4}{5} \rangle$.

3. \textbf{$D_{\mathbf{u}}f$:}
   \[ 1(-\frac{3}{5}) + 6(\frac{4}{5}) = \frac{-3+24}{5} = \frac{21}{5} \]

\textbf{Final Answer:}
$\nabla f = \langle 1, 6 \rangle$, $\mathbf{u} = \langle -0.6, 0.8 \rangle$, $D_{\mathbf{u}}f = \boxed{4.2}$

\hrulefill

\subsection*{Problem 17}
\textbf{$f(x, y, z) = xy^2z^3$, $P(3, 1, 1)$, $Q(0, -5, 7)$.}

1. \textbf{Gradient at P:}
   $f_x = y^2z^3 = 1$, $f_y = 2xyz^3 = 6$, $f_z = 3xy^2z^2 = 9$.
   $\nabla f = \langle 1, 6, 9 \rangle$.

2. \textbf{Direction P to Q:}
   $\mathbf{v} = \langle -3, -6, 6 \rangle$. $|\mathbf{v}| = \sqrt{9+36+36} = \sqrt{81} = 9$.
   $\mathbf{u} = \langle -3/9, -6/9, 6/9 \rangle = \langle -1/3, -2/3, 2/3 \rangle$.

3. \textbf{Dot Product:}
   \[ 1(-1/3) + 6(-2/3) + 9(2/3) = -1/3 - 4 + 6 = 2 - 1/3 = 5/3 \]

\textbf{Final Answer:}
\[ \boxed{\frac{5}{3}} \]

\hrulefill

\subsection*{Problem 18}
\textbf{Find max rate of change of $f(x, y) = 7xy^2$ at $(3, -7)$.}

Gradient $\nabla f = \langle 7y^2, 14xy \rangle$.
At $(3, -7)$: $\langle 7(49), 14(3)(-7) \rangle = \langle 343, -294 \rangle$.
Max rate = $|\nabla f| = \sqrt{343^2 + (-294)^2}$.
Factoring out $49$: $\langle 49(7), 49(-6) \rangle$.
Mag = $49\sqrt{7^2 + (-6)^2} = 49\sqrt{49+36} = 49\sqrt{85}$.
Direction: $\langle 343, -294 \rangle$ or normalized.

\textbf{Final Answer:}
Rate: $\boxed{49\sqrt{85}}$
Direction: $\boxed{\langle 343, -294 \rangle}$

\hrulefill

\subsection*{Problem 19}
\textbf{Max rate of $f(s, t) = \frac{s}{s^2+t^2}$ at $(-6, 6)$.}

Gradient Quotient Rule. Denom $D = s^2+t^2 = 36+36=72$.
$f_s = \frac{1(D) - s(2s)}{D^2} = \frac{t^2-s^2}{D^2}$.
$f_t = \frac{-s(2t)}{D^2} = \frac{-2st}{D^2}$.
At $(-6, 6)$:
$t^2 - s^2 = 36 - 36 = 0$. So $f_s = 0$.
$-2st = -2(-6)(6) = 72$. So $f_t = 72 / 72^2 = 1/72$.
$\nabla f = \langle 0, 1/72 \rangle$.
Max rate: $1/72$. Direction $\langle 0, 1/72 \rangle$.

\textbf{Final Answer:}
Rate: $\boxed{\frac{1}{72}}$, Direction: $\boxed{\langle 0, \frac{1}{72} \rangle}$ or $\mathbf{j}$.

\hrulefill

\subsection*{Problem 20}
\textbf{Max rate of $f(x, y) = 2\sin(xy)$ at $(0, 3)$.}

$\nabla f = \langle 2y \cos(xy), 2x \cos(xy) \rangle$.
At $(0, 3)$: $\langle 6 \cos(0), 0 \cos(0) \rangle = \langle 6, 0 \rangle$.
Max rate = 6. Direction $\langle 6, 0 \rangle$ or $\mathbf{i}$.

\textbf{Final Answer:}
Rate: $\boxed{6}$, Direction: $\boxed{\langle 6, 0 \rangle}$

\hrulefill

\subsection*{Problem 21}
\textbf{Max rate of $f(x, y, z) = x \ln(yz)$ at $(2, 5, 1/5)$.}

$yz = 1$. $\ln(1)=0$.
$f_x = \ln(yz) = 0$.
$f_y = x(1/yz)(z) = x/y = 2/5$.
$f_z = x(1/yz)(y) = x/z = 2/(1/5) = 10$.
$\nabla f = \langle 0, 0.4, 10 \rangle$.
Rate = $\sqrt{0 + 0.16 + 100} = \sqrt{100.16}$.
Direction: $\langle 0, 2/5, 10 \rangle$.

\textbf{Final Answer:}
Rate: $\boxed{\sqrt{100.16}}$, Direction: $\boxed{\langle 0, 2/5, 10 \rangle}$

\hrulefill

\subsection*{Problem 22}
\textbf{Find points where fastest change direction of $f(x, y) = x^2 + y^2 - 10x - 12y$ is $\mathbf{i} + \mathbf{j}$.}

$\nabla f = \langle 2x - 10, 2y - 12 \rangle$.
For the direction to be $\mathbf{i} + \mathbf{j}$, $\nabla f$ must be a positive scalar multiple of $\langle 1, 1 \rangle$.
$2x - 10 = k$ and $2y - 12 = k$ with $k > 0$.
$2x - 10 = 2y - 12 \implies 2y = 2x + 2 \implies y = x + 1$.
Condition $k > 0 \implies 2x - 10 > 0 \implies x > 5$.

\textbf{Final Answer:}
All points on the line $\boxed{y = x + 1}$ where $\boxed{x > 5}$.

\hrulefill

\subsection*{Problem 23}
\textbf{$T = 100e^{-x^2-3y^2-7z^2}$ at $P(4, -1, 4)$.}

Exponent $E = -16 - 3 - 112 = -131$.
$\nabla T = T \cdot \langle -2x, -6y, -14z \rangle$.
At P: $100e^{-131} \langle -8, 6, -56 \rangle$.

(a) Towards $(6, -2, 6)$. $\mathbf{v} = \langle 2, -1, 2 \rangle$. $|\mathbf{v}|=3$. $\mathbf{u}=\mathbf{v}/3$.
$D_{\mathbf{u}}T = \frac{100e^{-131}}{3} [-8(2) + 6(-1) - 56(2)] = \frac{100e^{-131}}{3} [-16 - 6 - 112] = \frac{100e^{-131}}{3}(-134)$.

(b) Direction of fastest increase: $\nabla T = \langle -8, 6, -56 \rangle$.

(c) Max rate: $|\nabla T| = 100e^{-131} \sqrt{64 + 36 + 3136} = 100e^{-131} \sqrt{3236}$.

\textbf{Final Answer:}
(a) $\boxed{-\frac{13400}{3}e^{-131}}$, (b) $\boxed{\langle -8, 6, -56 \rangle}$, (c) $\boxed{100e^{-131}\sqrt{3236}}$

\hrulefill

\subsection*{Problem 24}
\textbf{$V = 5x^2 - 3xy + xyz$ at $P(6, 2, 5)$. Direction $\mathbf{i}+\mathbf{j}-\mathbf{k}$.}

$\nabla V = \langle 10x - 3y + yz, -3x + xz, xy \rangle$.
At $(6, 2, 5)$:
$V_x = 60 - 6 + 10 = 64$.
$V_y = -18 + 30 = 12$.
$V_z = 12$.
$\mathbf{u} = \frac{1}{\sqrt{3}}\langle 1, 1, -1 \rangle$.
(a) $D_{\mathbf{u}}V = \frac{1}{\sqrt{3}}(64 + 12 - 12) = \frac{64}{\sqrt{3}}$.
(b) Direction: $\langle 64, 12, 12 \rangle$.
(c) Max Rate: $\sqrt{64^2 + 12^2 + 12^2} = \sqrt{4096 + 144 + 144} = \sqrt{4384}$.

\textbf{Final Answer:}
(a) $\boxed{\frac{64}{\sqrt{3}}}$, (b) $\boxed{\langle 64, 12, 12 \rangle}$, (c) $\boxed{\sqrt{4384}}$

\hrulefill

\subsection*{Problem 25}
\textbf{Second Directional Derivative $D_{\mathbf{u}}^2 f$. $f = x^3 + 5x^2y + y^3$. $\mathbf{u} = \langle 5/13, 12/13 \rangle$.}

Formula: $D_{\mathbf{u}}^2 f = f_{xx}a^2 + 2f_{xy}ab + f_{yy}b^2$.
$f_{xx} = 6x + 10y$. $f_{xy} = 10x$. $f_{yy} = 6y$.
At $(1, 3)$:
$f_{xx} = 6 + 30 = 36$.
$f_{xy} = 10$.
$f_{yy} = 18$.
$a = 5/13, b = 12/13$.
Value = $\frac{1}{169} [ 36(25) + 2(10)(60) + 18(144) ] = \frac{900 + 1200 + 2592}{169} = \frac{4692}{169}$.

\textbf{Final Answer:}
\[ \boxed{\frac{4692}{169}} \approx 27.76 \]

\hrulefill

\subsection*{Problems 26-28: Tangent Planes}
\textbf{Formula:} $\nabla F(P) \cdot \langle x-x_0, y-y_0, z-z_0 \rangle = 0$.

\textbf{26.} $F = 2(x-5)^2 + (y-7)^2 + (z-5)^2 = 10$ at $(6, 9, 7)$.
$\nabla F = \langle 4(x-5), 2(y-7), 2(z-5) \rangle$.
At $(6, 9, 7)$: $\langle 4(1), 2(2), 2(2) \rangle = \langle 4, 4, 4 \rangle$.
Tangent Plane: $4(x-6) + 4(y-9) + 4(z-7) = 0 \implies x+y+z = 22$.
Normal Line: $x=6+4t, y=9+4t, z=7+4t$.

\textbf{27.} Same logic. $\nabla F = \langle 4(x-6), 2(y-6), 2(z-5) \rangle$. At $(7, 8, 7)$: $\langle 4, 4, 4 \rangle$.
Same normal vector.

\textbf{28.} $x+y+z - 6e^{xyz} = 0$.
$\nabla F = \langle 1 - 6yz e^{xyz}, 1 - 6xz e^{xyz}, 1 - 6xy e^{xyz} \rangle$.
At $(0, 0, 6)$: $\langle 1, 1, 1 \rangle$.
Plane: $1(x-0) + 1(y-0) + 1(z-6) = 0 \implies x+y+z=6$.
Line: $x=t, y=t, z=6+t$.

\hrulefill

\subsection*{Problem 29}
\textbf{Tangent to Level Curve.}
$g(x, y) = x^2 + y^2 - 6x$. $\nabla g = \langle 2x-6, 2y \rangle$.
At $(2, 6)$: $\nabla g = \langle -2, 12 \rangle$.
Tangent line is perpendicular to $\langle -2, 12 \rangle$. Equation: $-2(x-2) + 12(y-6) = 0 \implies -x + 6y = 34$.
Sketch: The level curve is a circle centered at $(3, 0)$. Gradient points outward/inward depending on values. At $(2, 6)$, gradient is $\langle -2, 12 \rangle$ (Left and Up).
Correct graph is Top-Left in provided options.

\newpage

\section{Analysis and "Cheatsheet"}

\subsection{Problem Types and Analysis}
\begin{enumerate}
    \item \textbf{Angle Problems (Problems 1-5):} Given $\theta$. Use $\mathbf{u} = \langle \cos\theta, \sin\theta \rangle$. \textit{Trap:} Ensure $\theta$ is in radians.
    \item \textbf{Gradient Evaluation (Problems 6-9, 29):} Just finding $\nabla f$. Straightforward partial differentiation.
    \item \textbf{Non-Unit Vector Direction (Problems 10-14):} Given $\mathbf{v}$, you must calculate $\mathbf{u} = \mathbf{v}/|\mathbf{v}|$. This is the most common student error.
    \item \textbf{Two-Point Direction (Problems 15-17):} Direction is vector $\vec{PQ} = Q - P$. Normalize this vector.
    \item \textbf{Max Rate of Change (Problems 18-21, 23, 24):} Magnitude $|\nabla f|$. Direction is $\nabla f$ itself.
    \item \textbf{Tangent Planes/Lines (Problems 26-28):} Use gradient as normal vector $\mathbf{n}$. Plane equation is $\mathbf{n} \cdot (\mathbf{r} - \mathbf{r}_0) = 0$.
\end{enumerate}

\subsection{Cheatsheet for Success}
\begin{itemize}
    \item \textbf{Is it a unit vector?} Check if $a^2 + b^2 = 1$. If not, DIVIDE by the length!
    \item \textbf{Gradient Geometry:} Gradient is perpendicular to level curves/surfaces.
    \item \textbf{Max Rate:} Simply $|\nabla f|$.
    \item \textbf{Direction of Max Rate:} simply $\nabla f$ (as a vector).
    \item \textbf{Zero Change:} Occurs in direction perpendicular to gradient.
\end{itemize}

\section{Real-World Application: Financial Derivatives}
In Quantitative Finance, the price of an option $V$ depends on the stock price $S$, volatility $\sigma$, and time $t$.
A "Greek" like Delta is just a partial derivative $\Delta = \frac{\partial V}{\partial S}$.
However, markets move in complex ways. A \textbf{Directional Derivative} allows a risk manager to estimate the change in portfolio value if, for example, the stock price drops AND volatility spikes simultaneously (a common crash scenario).
\[ D_{\mathbf{u}}V \approx \frac{\partial V}{\partial S} \Delta S + \frac{\partial V}{\partial \sigma} \Delta \sigma \]
This models the "slope" of the value function in a specific crash-scenario direction.

\section{Advanced Diagnostic: Find the Flaw}

\textbf{Problem 1:} Find $D_{\mathbf{u}}f$ for $f=x^2$ in direction $\mathbf{v}=\langle 3, 4 \rangle$ at $(1, 0)$.
\textit{Flawed Solution:} $\nabla f = \langle 2x, 0 \rangle = \langle 2, 0 \rangle$. $D_{\mathbf{u}}f = \langle 2, 0 \rangle \cdot \langle 3, 4 \rangle = 6$.
\textit{Error:} Failed to normalize $\mathbf{v}$.
\textit{Correction:} $\mathbf{u} = \langle 3/5, 4/5 \rangle$. Answer: $6/5$.

\textbf{Problem 2:} Find direction of max change.
\textit{Flawed Solution:} Max change is $|\nabla f| = 10$.
\textit{Error:} Confused "direction" with "rate".
\textit{Correction:} Direction is the vector $\nabla f$, not its magnitude.

\textbf{Problem 3:} Tangent plane to $z = x^2+y^2$ at $(1,1,2)$.
\textit{Flawed Solution:} Gradient of $x^2+y^2$ is $\langle 2x, 2y \rangle = \langle 2, 2 \rangle$. Plane: $2(x-1)+2(y-1)=0$.
\textit{Error:} Used 2D gradient for 3D surface. Must define $F(x, y, z) = x^2+y^2-z=0$.
\textit{Correction:} $\nabla F = \langle 2, 2, -1 \rangle$. Plane: $2(x-1)+2(y-1)-1(z-2)=0$.

\textbf{Problem 4:} $D_{\mathbf{u}}f$ with angle $\theta=30^\circ$.
\textit{Flawed Solution:} $\mathbf{u} = \langle \sin(30), \cos(30) \rangle$.
\textit{Error:} Swapped sine and cosine. $x$ goes with cosine.
\textit{Correction:} $\mathbf{u} = \langle \cos(30), \sin(30) \rangle$.

\textbf{Problem 5:} Max rate of change is 0.
\textit{Flawed Solution:} $\nabla f = \mathbf{0}$. So max rate is undefined.
\textit{Error:} Max rate is 0. It is defined, just zero. (Critical point).

\end{document}